% Sección 7.6: Ejercicios

\subsection*{SECCIÓN 7.6 \quad Ejercicios}

\noindent\textbf{Evalúe las siguientes integrales.}
\begin{enumerate}
    \subimport{ejercicio_01/}{ejercicio.tex}
\subimport{ejercicio_02/}{ejercicio.tex}
\subimport{ejercicio_03/}{ejercicio.tex}
\subimport{ejercicio_04/}{ejercicio.tex}
\subimport{ejercicio_05/}{ejercicio.tex}
\subimport{ejercicio_06/}{ejercicio.tex}
\subimport{ejercicio_07/}{ejercicio.tex}
\subimport{ejercicio_08/}{ejercicio.tex}
\subimport{ejercicio_09/}{ejercicio.tex}
\subimport{ejercicio_10/}{ejercicio.tex}
\subimport{ejercicio_11/}{ejercicio.tex}
\subimport{ejercicio_12/}{ejercicio.tex}
\subimport{ejercicio_13/}{ejercicio.tex}
\subimport{ejercicio_14/}{ejercicio.tex}
\subimport{ejercicio_15/}{ejercicio.tex}
\subimport{ejercicio_16/}{ejercicio.tex}
\subimport{ejercicio_17/}{ejercicio.tex}
\subimport{ejercicio_18/}{ejercicio.tex}
\subimport{ejercicio_19/}{ejercicio.tex}
\subimport{ejercicio_20/}{ejercicio.tex}
\subimport{ejercicio_21/}{ejercicio.tex}
\subimport{ejercicio_22/}{ejercicio.tex}
\subimport{ejercicio_23/}{ejercicio.tex}
\subimport{ejercicio_24/}{ejercicio.tex}
\subimport{ejercicio_25/}{ejercicio.tex}
\subimport{ejercicio_26/}{ejercicio.tex}
\subimport{ejercicio_27/}{ejercicio.tex}
\subimport{ejercicio_28/}{ejercicio.tex}
\subimport{ejercicio_29/}{ejercicio.tex}
\subimport{ejercicio_30/}{ejercicio.tex}
\subimport{ejercicio_31/}{ejercicio.tex}
\subimport{ejercicio_32/}{ejercicio.tex}
\subimport{ejercicio_33/}{ejercicio.tex}
\subimport{ejercicio_34/}{ejercicio.tex}
\subimport{ejercicio_35/}{ejercicio.tex}
\subimport{ejercicio_36/}{ejercicio.tex}
\subimport{ejercicio_37/}{ejercicio.tex}
\subimport{ejercicio_38/}{ejercicio.tex}
\subimport{ejercicio_39/}{ejercicio.tex}
\subimport{ejercicio_40/}{ejercicio.tex}
\subimport{ejercicio_41/}{ejercicio.tex}
\subimport{ejercicio_42/}{ejercicio.tex}
\end{enumerate}

\vspace{0.5cm}
\begin{enumerate}
    \setcounter{enumi}{42}
    \subimport{ejercicio_43/}{ejercicio.tex}
\subimport{ejercicio_44/}{ejercicio.tex}
\subimport{ejercicio_45/}{ejercicio.tex}
\subimport{ejercicio_46/}{ejercicio.tex}
\subimport{ejercicio_47/}{ejercicio.tex}
\subimport{ejercicio_48/}{ejercicio.tex}
\subimport{ejercicio_49/}{ejercicio.tex}
\subimport{ejercicio_50/}{ejercicio.tex}
\subimport{ejercicio_51/}{ejercicio.tex}
\subimport{ejercicio_52/}{ejercicio.tex}
\subimport{ejercicio_53/}{ejercicio.tex}
\subimport{ejercicio_54/}{ejercicio.tex}
\subimport{ejercicio_55/}{ejercicio.tex}
\subimport{ejercicio_56/}{ejercicio.tex}
\subimport{ejercicio_57/}{ejercicio.tex}
\subimport{ejercicio_58/}{ejercicio.tex}
\subimport{ejercicio_59/}{ejercicio.tex}
\subimport{ejercicio_60/}{ejercicio.tex}
\subimport{ejercicio_61/}{ejercicio.tex}
\subimport{ejercicio_62/}{ejercicio.tex}
\subimport{ejercicio_63/}{ejercicio.tex}
\subimport{ejercicio_64/}{ejercicio.tex}
\subimport{ejercicio_65/}{ejercicio.tex}
\subimport{ejercicio_66/}{ejercicio.tex}
\subimport{ejercicio_67/}{ejercicio.tex}
\subimport{ejercicio_68/}{ejercicio.tex}
\subimport{ejercicio_69/}{ejercicio.tex}
\subimport{ejercicio_70/}{ejercicio.tex}
\subimport{ejercicio_71/}{ejercicio.tex}
\subimport{ejercicio_72/}{ejercicio.tex}
\subimport{ejercicio_73/}{ejercicio.tex}
\subimport{ejercicio_74/}{ejercicio.tex}
\subimport{ejercicio_75/}{ejercicio.tex}
\subimport{ejercicio_76/}{ejercicio.tex}
\subimport{ejercicio_77/}{ejercicio.tex}
\subimport{ejercicio_78/}{ejercicio.tex}
\subimport{ejercicio_79/}{ejercicio.tex}
\subimport{ejercicio_80/}{ejercicio.tex}
\end{enumerate}

\vspace{0.8cm}
\subsection*{SECCIÓN 7.7 \quad Uso de Tablas de Integrales}

El científico o el ingeniero a menudo emplean tablas de integrales para integrar funciones complicadas. En el apéndice G se proporciona una tabla relativamente breve de 112 integrales. Tablas más extensas están disponibles en las \textit{CRC Mathematical Tables} (463 integrales) o en Gradshteyn and Ryzhik's \textit{Table of Integrals, Series and Products} (Nueva York: Academic Press, 1979), que contiene centenares de páginas de integrales. Se debe recordar, sin embargo, que las integrales no aparecen a menudo exactamente en la forma registrada en una tabla. Normalmente uno de los métodos de este capítulo, como el de sustitución o de integración por partes, se requiere para transformar una integral dada en una de las formas contenidas en la tabla.
